\section{Results}
\label{sec:results}


\subsection{Context-length ablation}
Extending the text encoder context length consistently improves retrieval performance and training efficiency.
Table~\ref{tab:retrieval_all} shows the zero-shot image-to-text and text-to-image recall at $k$ in the long context benchmarks. 
Panel A shows that longer context improves retrieval across recall levels for both CXR and PMC.
On CXR, gains are steady but relatively modest. 
PMC benefits most, especially at stricter thresholds, highlighting the value of long contexts for text-heavy tasks.
In addition, we observe that pretraining CLIP with longer context lengths accelerates convergence (appendix section~\ref{apd:converge}), indicating that context extension improves not only downstream retrieval but also training efficiency.

% \begin{figure*}[htbp]
%   \centering
%  % \fbox{\rule{0pt}{200pt}\rule{150pt}{0pt}} % height=100pt, width=150pt
%   \includegraphics[width=0.8\linewidth]{}
%   \caption{
%   % (A) BIOMEDICA-6M caption token length distribution 
%   Zero-shot accuracy as a function of text context length. A) Zero-shot image classification. B) Zero-shot image-text and text-image retrievel.}
%   \label{fig: motivation-and-example}
% \end{figure*}



\subsection{Benchmarking against baselines}
BMC-LongCLIP outperforms prior biomedical VLMs on both long-text benchmarks. Table~\ref{tab:retrieval_all} Panel B and Table~\ref{tab:zeroshot} compare BMC-LongCLIP with existing biomedical VLMs, including PMC-CLIP, BiomedCLIP, MedSigLIP, and BMC-CLIP. We exclude MedSigLIP from the CXR benchmark comparison, as it was trained on the same MIMIC-CXR image–report pairs used to construct our benchmark.


% On the CXR benchmark, baselines achieve $<$6\% Recall@10, while BMC-LongCLIP variants reach 10–14\%, a $>$2× improvement. On PMC, BMC-LongCLIP reaches 89–95\% R@10, matching/surpassing BiomedCLIP and far outperforming MedSigLIP; gains are largest at stricter thresholds.

On the CXR benchmark, baselines achieve $<$6\% Recall@10, while BMC-LongCLIP variants reach 10–14\%, a more than two-fold improvement. 
On PMC, BMC-LongCLIP achieves 89–95\% R@10, performing on par with or slightly better than BiomedCLIP (91–94\%) and outperforming MedSigLIP (46–60\%). At the stricter R@1 threshold, BMC-LongCLIP attains 69–81\%, surpassing BiomedCLIP (69–73\%) and outperforming MedSigLIP (20–31\%).

% On PMC, BMC-LongCLIP attains 89–95\% R@10, matching or slightly surpassing BiomedCLIP (91–94\%) and outperforming MedSigLIP (46–60\%), with the largest relative gains at stricter thresholds. For R@1, BMC-LongCLIP reaches 69–81\%, exceeding both BiomedCLIP (64–72\%) and far surpassing MedSigLIP (21–36\%)




% On the CXR benchmark, baseline models perform poorly, with Recall@10 values below 6\%. In contrast, BMC-LongCLIP-512 variants achieve Recall@10 between 10.3--12.2\% (T2I) and 9.8--14.5\% (I2T). Across all BMC-LongCLIP variants, this represents more than a two-fold improvement over the strongest baseline.

% On the PMC benchmark, BMC-LongCLIP variants achieve Recall@10 scores ranging from 89.3\% to 95.1\% (T2I) and 89.9\% to 93.8\% (I2T). These results closely match or slightly surpass BiomedCLIP (91.1 and 93.7\%) and substantially outperform MedSigLIP (45.5 and 60.7\%). While we highlight R@10 for clarity, the same trend holds across all recall levels, with the largest relative gains observed at stricter thresholds such as R@1.

Beyond retrieval, BMC-LongCLIP also improves zero-shot classification accuracy across six biomedical domains (Table~\ref{tab:zeroshot}). 
While BiomedCLIP and MedSigLIP achieve average accuracies of 41.9\% and 36.6\%, respectively, BMC-LongCLIP (8K) attains 50.2\%, the best overall performance. 
These results highlight that extending context length not only benefits long-text retrieval but also provides improvements in classification tasks.

\subsection{Effect of batch size and long-caption training} Long-context models benefit from enriched captions, while larger batch sizes yield mixed results. Doubling batch size with BMC-LongCLIP (16K) underperforms in microscopy and dermatology. This suggests that long context windows combined with large batch sizes may not uniformly translate into performance gains across domains. In contrast, BMC-LongCLIP+ (16K) with BIOMEDICA-LongCAP data recovers this drop and matches or exceeds the 8K model. These results indicate that long-context modeling is most effective when paired with sufficient long-caption supervision, though performance in microscopy remains comparatively weak and warrants further investigation.
\noindent \textbf{Overall.} 
Across all experiments, BMC-LongCLIP outperforms prior biomedical VLMs in long-text retrieval and provides competitive advantages in zero-shot classification.